\section{Aircraft specification and composition}
\label{sec:configuration}

Simulation instances are built from an \code{AircraftSpec} (\code{src/sim/Aircraft/Spec.jl}) typically loaded from TOML via \code{Aircraft.load\_spec} (\code{src/sim/Aircraft/TOMLIO.jl}). The spec layer exists to separate:
\begin{itemize}[leftmargin=*,itemsep=2pt]
  \item continuous plant integration choices (integrator, contact model),
  \item environment and timeline definition,
  \item aircraft component instances (airframe geometry, motors, batteries), and
  \item runtime boundary configuration (PX4 library path, uORB config, logging/recording).
\end{itemize}

\subsection{TOML specs and inheritance}

The repository ships TOML-backed default specs for common workflows (e.g. Iris). Specs support an \code{extends=[...]} field which deep-merges one or more base specs to form the final configuration (\code{_load\_toml\_with\_extends}). The merge rules are:
\begin{itemize}[leftmargin=*,itemsep=2pt]
  \item tables merge recursively,
  \item arrays replace (they do \emph{not} append),
  \item scalars replace.
\end{itemize}
Cycles in the \code{extends} chain are detected and rejected.

\paragraph{Strict schema checking.}
When \code{strict=true} (the default in all shipped workflows), each TOML table is guarded by an explicit allowed-key set (\code{_known\_keys!}). Unknown keys are treated as errors rather than being ignored.

\paragraph{Path resolution.}
All file paths in the spec (e.g. \code{px4.mission\_path}, \code{px4.libpath}, battery asset CSVs) are resolved relative to the directory of the root spec file (\code{_resolve\_path}) with \code{\textasciitilde} expansion. Some fields require the file to exist (e.g. \code{extends} items), others allow non-existing outputs (e.g. log/record output paths).

\paragraph{Microsecond quantization.}
Timeline intervals in seconds are converted to integer microseconds with a strict exactness check (\cref{sec:runtime-timebase}). If you specify a cadence such as \(0.004\) seconds, it must be exactly representable as an integer number of microseconds.

\subsection{Assembly into an engine}

The build entrypoint constructs an \emph{aircraft instance} containing:
\begin{itemize}[leftmargin=*,itemsep=2pt]
  \item a fully constructed hybrid timeline (autopilot/wind/log/scenario axes merged into a single event axis),
  \item an initial plant state \(x_0\),
  \item the plant RHS functor \(f(t,x,u)\) (typically \code{CoupledMultirotorModel}),
  \item an integrator instance, and
  \item a set of runtime sources (scenario, wind, estimator, autopilot) selected by run mode.
\end{itemize}
The runtime engine then wires these into \code{Sim.Runtime.Engine}.

\subsection{Schema overview (TOML \texorpdfstring{$\to$}{->} \texorpdfstring{\code{AircraftSpec}}{AircraftSpec})}
\label{sec:config-schema}

The top-level TOML keys map to strongly-typed sub-structs inside \code{AircraftSpec} via \code{spec\_from\_toml\_dict}. \cref{tab:schema-overview} summarizes the main blocks; each block has additional substructure that is validated during parsing.

\begin{table}[t]
  \centering
  \small
  \begin{tabularx}{\linewidth}{@{}l l X@{}}
    \toprule
    TOML block & Julia type & Notes (defaults / build-time behavior) \\
    \midrule
    \code{schema\_version} & \code{Int} & Must be \code{1}. \\
    \code{aircraft} & name + seed & \code{seed} drives deterministic RNGs (wind uses \code{seed}, estimator uses \code{seed+1}). \\
    \code{home} & \code{Autopilots.HomeLocation} & Used for NED$\leftrightarrow$LLA injection into PX4. \\
    \code{px4} & \code{PX4Spec} & \code{libpath} required for \code{:live}/\code{:record}; \code{uorb} contract defines pubs/subs. \\
    \code{timeline} & \code{TimelineSpec} & Defines \code{t\_end} and axis cadences; converted to integer microseconds. \\
    \code{environment} & \code{EnvironmentSpec} & Live build uses selected wind model; replay build forces \code{NoWind}. \\
    \code{scenario} & \code{ScenarioSpec} & Default builder creates an \code{EventScenario} with \code{arm\_time} and \code{mission\_time}. \\
    \code{estimator} & \code{EstimatorSpec} & \code{kind=:none} or \code{:noisy\_delayed}; delay defaults to \(2\,dt_{ap}\) when unset. \\
    \code{plant} & \code{PlantSpec} & Selects integrator and contact model. \\
    \code{airframe} & \code{AirframeSpec} & Multirotor geometry + propulsor + rigid-body parameters. \\
    \code{actuation} & \code{ActuationSpec} & Maps PX4 channels to physical motors/servos; selects actuator dynamics models. \\
    \code{power} & \code{PowerSpec} & Battery models + bus topology and sharing policy. \\
    \code{sensors} & sensor list & Optional; currently GPS/rangefinder/radar stubs. \\
    \code{run} & run defaults & Optional defaults for \code{mode}, recording paths, and CSV logging. \\
    \bottomrule
  \end{tabularx}
  \caption{High-level TOML schema mapping into \code{AircraftSpec}. The code treats this schema as part of the simulator contract: unknown keys are rejected in strict mode.}
  \label{tab:schema-overview}
\end{table}


\subsection{Live vs replay builds}

The same plant model and engine are used in multiple modes:
\begin{itemize}[leftmargin=*,itemsep=2pt]
  \item \textbf{Live}: sources are stateful components driven by deterministic RNGs (wind, estimator), and PX4 is stepped through the lockstep ABI.
  \item \textbf{Record}: a live run with additional boundary-aligned stream capture.
  \item \textbf{Replay}: live sources are replaced by trace samplers (ZOH/sample-and-hold) so the plant is driven by recorded inputs.
\end{itemize}
In replay, the environment's wind model is replaced with \code{NoWind} so that all wind forcing comes from the recorded wind trace and the bus sample-and-hold input.

\subsection{Geometry-consistent PX4 configuration}

When enabled, the builder derives PX4 control allocator geometry parameters from the multirotor spec (\cref{sec:px4-bridge}). This prevents a common lockstep failure mode where the controller allocates using a different rotor layout than the plant uses.

\paragraph{Defaults.}
{\raggedright
Many \code{AircraftSpec} sub-structs have pragmatic defaults. For example, \code{PlantSpec.integrator=:RK45}. Environment defaults include \code{EnvironmentSpec.wind=:ou} and \code{EnvironmentSpec.gravity=:uniform}. The shipped Iris spec asset in \cref{app:iris} provides a concrete example.\par}
